% ============================================================
% CHAPTER 10: WEEKLY BREAKDOWN
% MCDA 5585 / 5586 Major Project Report
% Bhavik Kantilal Bhagat | A00494758 | Saint Mary's University
% ============================================================

\chapter{Weekly Breakdown}
\label{chap:weekly_breakdown}

My internship was structured across four delivery phases spanning 24~weeks.
The following narrative covers what I built each week, what blocked or
accelerated progress, and how priorities shifted in response to production
events and emerging technical constraints.

% ────────────────────────────────────────────────────────────
\section{Phase 1: Foundation \& CloudWatch Dashboards (Weeks 1--7, Mar~15 -- May~10)}
\label{sec:wb_phase1}
% ────────────────────────────────────────────────────────────

\subsection{Weeks 1--4: Onboarding and First Deliverables (Mar 15 -- Apr 11)}

The engagement opened in part-time mode with approximately 10--20 hours per
week. The first priority was access provisioning: AWS console, JIRA,
Confluence, and CodeCommit access requests were filed on Day~1 and cleared
over the following two weeks. The initial period was used productively for
architecture documentation review on Confluence --- building a working mental
model of the six platforms (\texttt{rpacfs1}, \texttt{ia-bpo},
\texttt{automationhub}, \texttt{meridian}, \texttt{citcoworks}) and the
Lambda, ECS, and EC2 workloads distributed across them.

\textbf{Week 2} produced the first concrete deliverable: the RPACFS1
CloudWatch pipeline health dashboard (IISS-435), built against the
\texttt{rpacfs1-cais-pricing-extract-genai-infrastructure} repository. The
dashboard included Lambda function monitoring, Logs Insights error widgets,
and 262 automated tests.

\textbf{Week 5} (Apr 13--18) was the full training week. Four structured
knowledge-transfer sessions introduced the domain context required to design
meaningful observability coverage for the IA-IT-SRE Infrastructure:

\begin{itemize}
    \item \textbf{RPA Support Model (Horace)} --- the Blue Prism and UiPath
          operational model, the pool/queue architecture, PROD and UAT folder
          hierarchy (\texttt{COE $\to$ DataOps $\to$ Processes / Assets / Rules /
          Queues}), the support escalation model (service desk $\to$ Horace $\to$
          Kishor), VPL queue monitoring (5~PM~EST daily), and Corporate Actions
          template management (Mandatory, Voluntary, Ops Transfer).

    \item \textbf{CitcoWorks and Event Store (Bhanu Teja Murari)} --- the
          CitcoWorks Knowledge Graph (L1--L5 hierarchy backed by Amazon Neptune),
          the Event Routing Service (ERS) and Action Processor (AP) deployment
          model via \texttt{deployment\_specs.yaml}, the AuthMaster authentication
          service, and the structured Event Store logging library used across all
          IA processes (mandatory fields: \texttt{tenant\_id},
          \texttt{knowledge\_datetime}, area, \texttt{subject},
          type, source, code, message,
          \texttt{run\_id}, \texttt{process\_identifier}).

    \item \textbf{Meridian Platform and Kafka Architecture (Meridian team)} ---
          the Meridian event pipeline: ERS stores events in ODL/Redis Elasticache
          then pushes to MSK Kafka topics; the OpenSearch Consumer indexes events
          for AXI to query; six ECS clusters host Meridian services
          (\texttt{meridian-api-cluster}, \texttt{meridian-cluster-compute},
          \texttt{meridian-cluster-services}, \texttt{meridian-services-cluster},
          \texttt{meridian-ui-cluster}, \texttt{meridian-compute-cluster}); Kafka
          consumer lag on these topics is the primary leading indicator of event
          processing health. The Meridian Incident occurred on Apr~17 during
          this week, making the observability gap immediately concrete.

    \item \textbf{\AE xeo Treasury Platform} --- Citco's cloud-based wire
          approval and fund movement platform, integrating via COMAPI and the
          \AE xeo Reporter Server REST API; key workloads include Active P\&L
          reports, report generation, and account manager configurations used by
          MESO and the E-Binder reconciliation processes.
\end{itemize}

Also on Apr~13, Kishor reviewed an early Lambda dashboard prototype and provided
the requirements that shaped IISS-487: add three additional platform tag
dimensions (BPM, CitcoWorks, Meridian), improve filter logic for multi-tag
drill-down, and create a Confluence Lambda inventory page (IISS-461). By week
end, IISS-469 was underway and the \texttt{SEARCH()} visibility gap for
inactive Lambda functions had been identified and documented.

\textbf{Week 7} built the ECS (IISS-484), EKS (IISS-485), and EC2 (IISS-483)
monitoring dashboards. The EC2 gap analysis discovered the
\texttt{BudgetCode="CFS/RPA"} tag mismatch on \texttt{rpacfs1} instances,
introducing the AppName fallback pattern that became foundational for
the inventory API.

\subsection{Phase 1 Total: 100 hours --- Key Insight}

By the end of Phase~1, the Lambda inventory across all six platforms was complete
and the SEARCH() visibility gap had established the design constraint for the
next phase: CloudWatch Metric Insights SQL rather than subscription filters.

% ────────────────────────────────────────────────────────────
\section{Phase 2: Grafana AMG Intensive Build (Weeks 8--11, May~6 -- May~27)}
\label{sec:wb_phase2}
% ────────────────────────────────────────────────────────────

\textbf{Week 8} was the first RPA support rotation. IISS-487 was elevated
to Critical priority on May~6. Infrastructure deployment could not begin
until the following Monday due to the support shift, but the week was used
to scaffold the Grafana repository and plan the four-layer architecture.

\textbf{Week 9} (May 10--16) was the single most infrastructure-dense week
of the internship. May~11 saw four CloudFormation stacks deployed, the Grafana
AMG workspace provisioned with its 16-role IAM agent permission model, and
Container Insights enabled on five of the six Meridian clusters. Two days
later, on May~13, the unexpected Grafana 10.4→12.4 upgrade broke all SQL
panels and variable interpolation --- approximately one working day was
consumed migrating queries to the stable Grafana~12.4 model.

By \textbf{Week 11}, the core platform was functional: MSK consumer lag
alarm active, composite alarms configured, and the first six dashboard rows
complete. The detection window for a Meridian Kafka congestion event had
dropped from the 3.5-hour Apr~17 gap to a projected sub-30-second alarm fire.

\subsection{Phase 2 Total: 157 hours --- Key Milestone}

The MSK consumer lag alarm armed itself against live production traffic
on May~11, 2026 --- the inflection point where the observability gap
identified by the Meridian Incident was closed.

% ────────────────────────────────────────────────────────────
\section{Phase 3: Dashboard Completion, Hotfixes \& SQS (Weeks 12--15, May~28 -- Jun~18)}
\label{sec:wb_phase3}
% ────────────────────────────────────────────────────────────

\textbf{Week 12} was the second RPA support rotation. The MESO ECS
autoscale issue (IISS-538) was resolved alongside support duties.

\textbf{Week 13} was the most incident-dense week of the internship.
The CAIS Deadline \texttt{SpnegoError} (IISS-706) required 7 hours of
investigation and hotfix work, identifying the module-import credential
caching pattern as the root cause. In parallel, the IISS-507 inventory
API was scaffolded using the Kiro spec-driven workflow (2.5 hours), and
the SQS monitoring feature (IISS-482) was completed with 960 automated tests.

\textbf{Week 14} delivered the final dashboard at version v728, closing
IISS-487 at exactly 140 hours. Deep-link panels, DocumentDB and Redis
restructuring, and the main branch merge completed the primary deliverable.

\subsection{Phase 3 Total: 160 hours --- Key Milestone}

IISS-487 closed at v728 on Jun~18, 2026. The full platform --- 14 rows,
110+ panels, 38 alarms, 4 CloudFormation stacks --- was in production.

% ────────────────────────────────────────────────────────────
\section{Phase 4: Inventory API, X-Ray Research \& Report (Weeks 16--24, Jun~19 -- Aug~31)}
\label{sec:wb_phase4}
% ────────────────────────────────────────────────────────────

\textbf{Week 16} saw the MESO month-end stuck tasks incident (IISS-741,
Critical) requiring manual retrigger of INNOCAP/SAMLMOS fund processes. ECS
discoverers were completed concurrently, confirming the 22-cluster inventory
and implementing the two-stage AppName filter for the shared BudgetCode environments.

\textbf{Week 18--19} added the CloudWatch Logs observability layer to the
Grafana platform: 30 Logs Insights saved queries across 7 platforms, a
log group URL encoding utility, and Grafana Row~17 for log-derived panels.

\textbf{Week 21} addressed the ML-RPA-Status LLM Agent (IISS-824):
11 critical bugs fixed (OOM Lambda, Athena SSE, SSL, idempotent KB sync),
a 1,140-line monolithic function split into 8 modules, and 194 unit tests
written with 99\% coverage.

\textbf{Weeks 22--24} transitioned to report writing, knowledge transfer,
and handover documentation. The Mitacs BSI renewal for
Sep~1 -- Dec~31, 2026 was confirmed, providing continuity for
the X-Ray tracing and log-group discovery work deferred from this engagement.

\subsection{Phase 4 Total: 480 hours --- Key Milestone}

The 562-test inventory API reached 99\% coverage on Jul~14, 2026, completing
the third and final technical deliverable of the internship.

% ────────────────────────────────────────────────────────────
\section{Internship Hours Summary}
\label{sec:wb_hours}
% ────────────────────────────────────────────────────────────

\begin{table}[ht]
\centering
\caption{Internship Hours by Phase}
\label{tab:wb_hours}
\rowcolors{2}{tablerowA}{tablerowB}
\begin{tabularx}{\textwidth}{lp{1.2cm}Xp{1.5cm}}
\hline
\rowcolor{smubrand}\textcolor{white}{\textbf{Phase}} & \textcolor{white}{\textbf{Weeks}} & \textcolor{white}{\textbf{Primary Deliverables}} & \textcolor{white}{\textbf{Hours}} \\
\hline
Phase 1: Foundation \& CloudWatch & 1--7 & IISS-435, IISS-455/461/469, IISS-482--486 scaffolding & 100 \\
Phase 2: Grafana AMG Build & 8--11 & IISS-487 core (4 CFN stacks, Rows 1--6, alarms) & 157 \\
Phase 3: Dashboard Completion \& Hotfixes & 12--15 & IISS-487 v728, IISS-706, IISS-538, IISS-482 SQS & 160 \\
Phase 4: Inventory API, Research \& Report & 16--24 & IISS-507, IISS-741, IISS-824, X-Ray specs, report & 480 \\
\textbf{Total} & \textbf{24} & & \textbf{897} \\
\hline
\end{tabularx}
\end{table}

The internship ran from Mar~15 to Aug~31, 2026, logging
\textbf{897 hours} across 24 weeks. The remaining 3 hours to reach the
900-hour commitment are covered by the knowledge-transfer and handover
sessions in the final days of Aug and the first days of September.
Complete itemised work logs are provided in the appendix and in the
accompanying \texttt{WorkLogs.xlsx} submission file.